\section{Costruzioni che usano limiti}
\Todo{molti esempi disseminati per il capitolo 1 sono stati architettati per evitare di usare il linguaggio dei limiti;	ora diventano molto più snelli. Metto qui un po' di esempi senza pretesa di ordine (andranno spostati qua e là in ogni caso).}
\begin{example}
	\(\bbN\) come coequalizzatore in \(\ctCat\); più in generale, gruppo libero come coequalizzatore in \(\ctCat\).
\end{example}
\begin{example}
	Coequalizzatori riflessivi in categorie della forma \(\ctMod(\Omega,s)\) costruiscono i coprodotti (e.g., in monoidi, gruppi, anelli, \dots).
\end{example}
\begin{proposition}
	\Todo{Categorie filtrate e setacciate}
\end{proposition}
\begin{definition}
	\Todo{Il funtore dei sottoggetti, stavolta per davvero}
\end{definition}
\begin{example}
	Spiga di un (pre)fascio
\end{example}
\begin{remark}
	I colimiti filtrati sono importanti perché la relazione che genera il quoziente che definisce il coeq è già di equivalenza (grazie agli assiomi di cat filtrata, si apprezza particolarmente bene per gli ordini diretti...)
\end{remark}
\begin{example}
	Completamenti adici
\end{example}
\begin{example}
	Algebre iniziali, coalgebre terminali
\end{example}
\begin{examples}
	Esempi di algebre iniziali e coalgebre terminali.
\end{examples}
\begin{theorem}
	Teorema di Adamek
\end{theorem}
\subsection{Limiti e colimiti in categorie concrete}
\subsection{Colimiti in categorie algebriche}
\begin{esercizi}
	\item
	\item
	\item
	\item
	\item
\end{esercizi}
